\documentclass[10pt,a4paper]{amsart}
\usepackage[margin=19mm]{geometry}
\usepackage{amsmath,amssymb,mathtools}
\usepackage[scr=rsfs,cal=euler]{mathalfa}
\usepackage{xfrac}
\usepackage[unicode,hidelinks,bookmarks=false]{hyperref}
\newcommand{\ie}{i.e.\ }
\newcommand{\cf}{cf.\ }
\newcommand{\resp}{respectively\ }
\newcommand{\Subsection}{\subsection}
\providecommand{\nobreakdash}{\nobreak\hbox{-}}
\setlength{\emergencystretch}{2em}
\multlinegap0pt
\usepackage{bm}
\usepackage{bbm}
\usepackage{dsfont}
\usepackage{xspace}
\usepackage{cancel}
\usepackage{mathtools}
\usepackage{amsthm}
\makeatletter
\@ifundefined{Theo}{\newtheorem{Theo}{Theo}}{}
\@ifundefined{Lem}{\newtheorem{Lem}[Theo]{Lemma}}{}
\makeatother
\newcommand{\NN}{\mathbb{N}}
\newcommand{\RR}{\mathbb{R}}
\newcommand{\eps}{\varepsilon}
\title[Zeros of GKP sequences of polynomials]{Zeros of GKP sequences of polynomials}
\author{Antonio J. Dur\'an, Mario P\'erez, Juan L. Varona}
\date{Source: arXiv:2607.06578v1 (CC BY 4.0)}
\begin{document}
\maketitle

\noindent\emph{Source and licence.} Zeros of GKP sequences of polynomials. Version arXiv:2607.06578v1 (CC BY 4.0), distributed under the Creative Commons Attribution 4.0 International licence, as recorded in the arXiv licence metadata for this exact version and shown on the abstract page https://arxiv.org/abs/2607.06578v1. Full rights notice: licenses/arxiv-2607.06578-CC-BY-4.0.txt.\par\smallskip

\noindent\textit{Editorial scope.} Counted unit: the Lem whose source label is \texttt{lemaQn>d} in the article (source0.tex lines 1609--1639, source label \texttt{lemaQn>d}), delivered together with the article's own complete original proof of it (source0.tex lines 1641--1755, one \texttt{proof} environment of 115 lines). No mathematics is written, repaired or re-derived: statement and proof are the article's own text line for line. Required context retained uncounted: lemaG (lines 1421--1443); Gj (lines 1428--1442); Jmn (lines 1595--1600). Every definition, display and label of those lines is retained unchanged. Omitted from this package and not redistributed: 1--1420, 1443--1594, 1601--1608, 1640--1640, 1756--2087. Nothing inside a retained excerpt is omitted or altered: every retained body is delivered exactly as the article's lines are written, so no omission receipt of any kind accompanies this package. Citations: the retained excerpts carry no citation command at all, so no bibliography entry of the article is transcribed and no reference is redistributed. Reference adaptation: none was needed and none was made; every reference inside the delivered unit resolves inside it, so no unresolved reference or citation is produced. Printed slips kept literal: no printed word, symbol, hypothesis or number of the counted statement or of its proof was altered. Layout and class: the article's own class is \texttt{amsart} with packages including geometry, amsmath, amsthm, amsfonts, amssymb, hyperref; the wrapper is the lane's pinned \texttt{amsart} editorial wrapper at 10pt on A4, which restates the theorem-like environments the retained text needs and adds the article's own macro definitions that the retained text needs (\texttt{\textbackslash NN}, \texttt{\textbackslash RR}, \texttt{\textbackslash eps}), with extra wrapper packages (bm, bbm, dsfont, xspace, cancel, mathtools). The delivered numbering is the unit's own. Assets: no retained span contains a figure environment, a graphics-inclusion command, a PDF-page inclusion, a table or a TikZ picture; no image file is copied into this package, the package declares no asset dependency and no image file is redistributed with it. Licence: version arXiv:2607.06578v1 of this article is distributed under the Creative Commons Attribution 4.0 International licence, as recorded in the arXiv licence metadata for this exact version and shown on the abstract page https://arxiv.org/abs/2607.06578v1. A full rights notice is delivered at licenses/arxiv-2607.06578-CC-BY-4.0.txt. Characters: every retained excerpt is pure ASCII, so the delivered engine, XeLaTeX with its default Latin Modern text font, reports no missing character.
\begin{Lem}
\label{lemaG}
Let $G : \NN \cup \{0\} \to (0,+\infty)$ be a (positive) function such that
\[
  G(j-1) G(j+1) < G(j)^2
\]
for every $j \in \NN$. Then:
\begin{enumerate}
\item\label{Gj}
For each $m \in \NN$ and $j \in \NN \cup \{0\}$,
\[
  G(j) \left(\frac{G(m-1)}{G(m)}\right)^j
  \le \frac{G(m-1)^m}{G(m)^{m-1}}.
\]
Moreover, the left-hand side is strictly increasing in $j$ for $j \le m-1$ and strictly decreasing for $j \ge m$. The inequality is an equality for $j = m-1$ and $j = m$; otherwise, it is a strict inequality.

\item\label{limite0}
For each $m \in \NN$,
\[
  \lim_{j \to +\infty} G(j) \left(\frac{G(m-1)}{G(m)}\right)^j = 0.
\]
\end{enumerate}
\end{Lem}

\begin{equation}
\label{Jmn}
  J_{m,n} = \Bigg\{ x \in \RR : \Bigg| \frac{x}{-\frac{F(m-1)}{F(m)} 
   \left(\frac{G(m-1)}{G(m)}\right)^{n-K}} - 1 \Bigg| < \eps
  \Bigg\}.
\end{equation}

\begin{Lem}
\label{lemaQn>d}
Let
\begin{itemize}
\item $F : \NN \cup \{0\} \to \RR \setminus \{0\}$ be an arbitrary function,

\item $G : \NN \cup \{0\} \to (0,+\infty)$ be a (positive) function such that for every $j \in \NN$,
\[
  G(j-1) G(j+1) < G(j)^2,
\]

\item $K$ be a real number (not necessarily positive),

\item $\tau_j(n) \in \RR$, for each $j, n \in \NN \cup \{0\}$, $j \le n$, such that for each $j$,
\[
  \lim_{n \to +\infty} \tau_j(n) = 1.
\]
\end{itemize}
For each $n, N \in \NN \cup \{0\}$, let
\[
  Q_{n,N}(x) = \sum_{j=0}^N F(j) G(j)^{n-K} \tau_j(n) x^j.
\]
Given $\eps \in (0,1)$ and $N \in \NN$, there exist some $n_0$ and $\delta > 0$ such that for every $n \ge n_0$:
\begin{enumerate}
\item
The polynomial $Q_{n,N}$ has exactly one (simple) zero in each interval $J_{m,n}$ ($m=1,2,\dots,N$) defined by~\eqref{Jmn}.

\item
$|Q_{n,N}(x)| \ge \delta G(0)^{n-K}$, if $x \notin \bigcup_{m=1}^N J_{m,n}$.
\end{enumerate}
\end{Lem}

\begin{proof}
Let us fix $\eps \in (0,1)$ and $N \in \NN$; for now, take any~$n$. Let us evaluate $Q_{n,N}$ in both endpoints of an interval $J_{m,n}$, with $m=1,2,\dots,N$:
\begin{multline*}
  Q_{n,N}\left( -(1\pm \eps) \tfrac{F(m-1)}{F(m)} \left( \tfrac{G(m-1)}{G(m)} \right)^{n-K} \right)
  \\
  = \sum_{j=0}^N (-1)^j (1\pm \eps)^j F(j) \tfrac{F(m-1)^j}{F(m)^j}
  \tau_j(n) G(j)^{n-K} \left( \tfrac{G(m-1)}{G(m)} \right)^{j(n-K)}
\end{multline*}
(take apart the terms $j=m-1$ and $j=m$)
\begin{multline*}
  = (-1)^{m-1} (1\pm \eps)^{m-1} \tfrac{F(m-1)^m}{F(m)^{m-1}}
  \tau_{m-1}(n) \left( \tfrac{G(m-1)^m}{G(m)^{m-1}} \right)^{n-K}
  \\
  + (-1)^m (1\pm \eps)^m \tfrac{F(m-1)^m}{F(m)^{m-1}}
  \tau_m(n) \left( \tfrac{G(m-1)^m}{G(m)^{m-1}} \right)^{n-K}
  \\
  + \sum_{\substack{0 \le j \le N \\ j\ne m-1, m}} (-1)^j (1\pm \eps)^j F(j) \tfrac{F(m-1)^j}{F(m)^j}
  \tau_j(n) G(j)^{n-K} \left( \tfrac{G(m-1)}{G(m)} \right)^{j(n-K)}
%\end{multline*}
\\
%\begin{multline*}
  = (-1)^{m-1} (1\pm \eps)^{m-1} \tfrac{F(m-1)^m}{F(m)^{m-1}}
  \left( \tfrac{G(m-1)^m}{G(m)^{m-1}} \right)^{n-K}
  \big( \tau_{m-1}(n) - (1 \pm \eps) \tau_m(n) \big)
  \\
  + \sum_{\substack{0 \le j \le N \\ j\ne m-1, m}} (-1)^j (1\pm \eps)^j F(j) \tfrac{F(m-1)^j}{F(m)^j}
  \tau_j(n) \left( G(j) \left(\tfrac{G(m-1)}{G(m)}\right)^j \right)^{n-K}.
\end{multline*}
Dividing by $\left( \tfrac{G(m-1)^m}{G(m)^{m-1}} \right)^{n-K}$ we obtain that
\begin{multline*}
  \left( \tfrac{G(m)^{m-1}}{G(m-1)^m} \right)^{n-K} 
  Q_{n,N}\left( -(1\pm \eps) \tfrac{F(m-1)}{F(m)} \left( \tfrac{G(m-1)}{G(m)} \right)^{n-K} \right)
  \\
  = (-1)^{m-1} (1\pm \eps)^{m-1} \tfrac{F(m-1)^m}{F(m)^{m-1}}
  \big( \tau_{m-1}(n) - (1 \pm \eps) \tau_m(n) \big)
  \\
  + \sum_{\substack{0 \le j \le N \\ j\ne m-1, m}} (-1)^j 
  (1\pm \eps)^j F(j) \tfrac{F(m-1)^j}{F(m)^j}
  \tau_j(n) \left(
  \frac{G(j) \Big(\tfrac{G(m-1)}{G(m)}\Big)^j}
  {\tfrac{G(m-1)^m}{G(m)^{m-1}}}
  \right)^{n-K}.
\end{multline*}
Let us take now limit in $n$, taking into account that:
\begin{itemize}
\item For each $j$, $\lim\limits_{n\to +\infty} \tau_j(n) = 1$.

\item According to Lemma~\ref{lemaG}, part~\eqref{Gj}, $\frac{G(j) \left(\tfrac{G(m-1)}{G(m)}\right)^j}
  {\tfrac{G(m-1)^m}{G(m)^{m-1}}} < 1$ for $j \ne m-1, m$.
\end{itemize}
Therefore, for each $m \in \{1,2,\dots,N\}$,
\begin{multline}
\label{limQnM}
  \lim_{n\to +\infty} \left( \tfrac{G(m)^{m-1}}{G(m-1)^m} \right)^{n-K} 
  Q_{n,N}\left( -(1\pm \eps) \tfrac{F(m-1)}{F(m)} 
  \left( \tfrac{G(m-1)}{G(m)} \right)^{n-K} \right)
  \\
  = (-1)^{m-1} (1\pm \eps)^{m-1} \tfrac{F(m-1)^m}{F(m)^{m-1}}
  (\mp \eps).
\end{multline}
This implies that, given $m \in \{1,2,\dots,N\}$ and taking $n$ large enough, both values
\begin{equation}
\label{QImndistintosigno}
  Q_{n,N}\left( -(1\pm \eps) \tfrac{F(m-1)}{F(m)} \left( \tfrac{G(m-1)}{G(m)} \right)^{n-K} \right)
\end{equation}
have opposite sign, that is, $Q_{n,N}$ changes its sign at the endpoints of $J_{m,n}$, so it has at least one zero in this interval. For $n$ large enough, this happens in all the $J_{m,n}$, with $m=1,2,\dots,N$. Since $Q_{n,N}$ has degree $N$, these zeros are unique (and simple). This proves the first part of the lemma.

If all zeros of $Q_{n,N}$ are real and simple, then all zeros of its derivative $Q_{n,N}'$ are simple as well, and there is exactly one zero between any two consecutive zeros of $Q_{n,N}$. Between any two consecutive zeros, $|Q_{n,N}|$ increases up to some point, then it decreases. In particular, in any interval $[a,b]$ on which $Q_{n,N}$ does not vanishes, it is plain that
\[
  |Q_{n,N}(x)| \ge \min\{ |Q_{n,N}(a)|, |Q_{n,N}(b)|\},
  \quad x \in [a,b].
\]

Thus, if $n$ is so large that the values~\eqref{QImndistintosigno} have opposite sign for $m=1,2,\dots,N$, it follows that
\[
  |Q_{n,N}(x)| \ge \min_{m=1,2,\dots,N}\left\{
  \left|Q_{n,N}\left( -(1\pm \eps) \tfrac{F(m-1)}{F(m)} 
  \left( \tfrac{G(m-1)}{G(m)} \right)^{n-K} \right)\right|
  \right\}
\]
for $x \notin \bigcup_{m=1}^N J_{m,n}$. According to Lemma~\ref{lemaG}, part~\eqref{Gj},
we have that
\[
  1 \ge G(0) \frac{G(m)^{m-1}}{G(m-1)^m}.
\]
Therefore, keeping in mind the limit~\eqref{limQnM}, we conclude that if $n$ is large enough, then for every $x \notin \bigcup_{m=1}^N J_{m,n}$,
\begin{multline*}
  |Q_{n,N}(x)| \ge \min_{m=1,2,\dots,N}\Big\{
  \Big|Q_{n,N}\Big( -(1\pm \eps) \tfrac{F(m-1)}{F(m)} \Big( \tfrac{G(m-1)}{G(m)} \Big)^{n-K} \Big)\Big|
  \Big\}
  \\
  \ge G(0)^{n-K} \min_{m=1,2,\dots,N}\Big\{
  \Big( \tfrac{G(m)^{m-1}}{G(m-1)^m} \Big)^{n-K}
  \Big| Q_{n,N}\Big( -(1\pm \eps) \tfrac{F(m-1)}{F(m)} 
  \Big( \tfrac{G(m-1)}{G(m)} \Big)^{n-K} \Big)\Big|
  \Big\}
  \\
  \ge G(0)^{n-K} \min_{m=1,2,\dots,N}\Big\{
  \frac{1}{2} \Big|
  (-1)^{m-1} (1\pm \eps)^{m-1} \tfrac{F(m-1)^m}{F(m)^{m-1}}
  (\mp \eps)
  \Big|
  \Big\}
  \\
  = G(0)^{n-K} \min_{m=1,2,\dots,N} \Big\{ \frac{\eps}{2} (1 - \eps)^{m-1} 
  \Big|\tfrac{F(m-1)^m}{F(m)^{m-1}} \Big| \Big\}.
\end{multline*}
This proves the second part of the lemma, with
\[
  \delta = \min_{m=1,2,\dots,N} \left\{
  \frac{\eps}{2} (1 - \eps)^{m-1} \left|\tfrac{F(m-1)^m}{F(m)^{m-1}}\right|
  \right\}.
  \qedhere
\]
\end{proof}

\end{document}
